\ifdefined\gkver\else\def\gkver{student}\fi
\documentclass[\gkver]{gaokaozhenti}
\begin{document}

\chapter{2011年海南}

\begin{enumerate}
\item
%% number: 1
%% paperName: 2011年海南高考第 1 题 3 分
%% typeId: 1
%% type: 单选题
%% chapter: 电场
%% point: 电势能电势差电场力做功
%% method: 无
%% score: 3
%% degree: 650
%% duplicateId: 0
%% body:
关于静电场，下列说法正确的是\xzanswer{D}
\fourchoices[answer=D]
{电势等于零的物体一定不带电}
{电场强度为零的点，电势一定为零}
{同一电场线上的各点，电势一定相等}
{负电荷沿电场线方向移动时，电势能一定增加}

%% answer: D
%% memo:
\memoanswer{
考察电场和电势概念及其电场力做功与电势能的关系，选 D。
}

\item
%% number: 2
%% paperName: 2011年海南高考第 2 题 3 分
%% typeId: 1
%% type: 单选题
%% chapter: 电路
%% point: 动态分析
%% method: 无
%% score: 3
%% degree: 650
%% duplicateId: 0
%% body:
如图，$E$ 为内阻不能忽略的电池，$R_{1}$、$R_{2}$、$R_{3}$ 为定值电阻，$\mathrm{S}_{0}$、$\mathrm{S}$ 为开关，V 与 A 分别为电压表与电流表。初始时 $\mathrm{S}_{0}$ 与 $\mathrm{S}$ 均闭合，现将 $\mathrm{S}$ 断开，则\xzanswer{B}
\onepicture[w=4.2cm]{figs/02.png}
\fourchoices[answer=B]
{V 的读数变大，A 的读数变小}
{V 的读数变大，A 的读数变大}
{V 的读数变小，A 的读数变小}
{V 的读数变小，A 的读数变大}

%% answer: B
%% memo:
\memoanswer{
$\mathrm{S}$ 断开，相当于外电阻变大，总电流减小，故路端电压 $U=E-Ir$ 增大，V 的读数变大。$\mathrm{S}$ 断开，使原来 $R_{2}$ 与 $R_{3}$ 并联的电阻变为 $R_{3}$ 的电阻，电阻变大了，$R_{3}$ 两端的电压增大，电流增大，A 的读数变大。
}

\item
%% number: 3
%% paperName: 2011年海南高考第 3 题 3 分
%% typeId: 1
%% type: 单选题
%% chapter: 电场
%% point: 电荷守恒定律
%% method: 无
%% score: 3
%% degree: 650
%% duplicateId: 0
%% body:
三个相同的金属小球 1、2、3 分别置于绝缘支架上，各球之间的距离远大于小球的直径。球 1 的带电量为 $q$，球 2 的带电量为 $nq$，球 3 不带电且离球 1 和球 2 很远，此时球 1、2 之间作用力的大小为 $F$。现使球 3 先与球 2 接触，再与球 1 接触，然后将球 3 移至远处，此时 1、2 之间作用力的大小仍为 $F$，方向不变。由此可知\xzanswer{D}
\fourchoices[answer=D]
{$n=3$}
{$n=4$}
{$n=5$}
{$n=6$}

%% answer: D
%% memo:
\memoanswer{
设 1、2 间距离为 $R$，则 $F=\frac{nq^{2}}{R^{2}}$。球 3 与球 2 接触后，它们所带的电量均为 $\frac{nq}{2}$；再使球 3 与球 1 接触后，它们所带的电量均为 $\frac{(n+2)q}{4}$。最后球 1、2 间的作用力 $F=\frac{n(n+2)q^{2}}{8R^{2}}$，由上面两式得 $n=6$。
}

\item
%% number: 4
%% paperName: 2011年海南高考第 4 题 3 分
%% typeId: 1
%% type: 单选题
%% chapter: 力物体的平衡
%% point: 三力平衡
%% method: 无
%% score: 3
%% degree: 450
%% duplicateId: 0
%% body:
如图，墙上有两个钉子 a 和 b，它们的连线与水平方向的夹角为 $45^\circ$，两者的高度差为 $l$。一条不可伸长的轻质细绳一端固定于 a 点，另一端跨过光滑钉子 b 悬挂一质量为 $m_{1}$ 的重物。在绳子距 a 端 $l/2$ 的 c 点有一固定绳圈。若绳圈上悬挂质量为 $m_{2}$ 的钩码，平衡后绳的 ac 段正好水平，则重物和钩码的质量比 $\frac{m_{1}}{m_{2}}$ 为\xzanswer{C}
\onepicture[w=3.2cm]{figs/04.png}
\fourchoices[answer=C]
{$\sqrt{5}$}
{$2$}
{$\frac{\sqrt{5}}{2}$}
{$\sqrt{2}$}

%% answer: C
%% memo:
\memoanswer{
平衡后设绳的 BC 段与水平方向成 $\alpha$ 角，则 $\tan\alpha=2$，$\sin\alpha=\frac{2}{\sqrt{5}}$。对节点 C 分析三力平衡，在竖直方向上有 $m_{2}g=m_{1}g\sin\alpha$，得 $\frac{m_{1}}{m_{2}}=\frac{1}{\sin\alpha}=\frac{\sqrt{5}}{2}$，选 C。
}

\item
%% number: 5
%% paperName: 2011年海南高考第 5 题 3 分
%% typeId: 1
%% type: 单选题
%% chapter: 力物体的平衡
%% point: 物体系平衡
%% method: 整体与隔离
%% score: 3
%% degree: 450
%% duplicateId: 0
%% body:
如图，粗糙的水平地面上有一斜劈，斜劈上一物块正在沿斜面以速度 $v_{0}$ 匀速下滑，斜劈保持静止，则地面对斜劈的摩擦力\xzanswer{A}
\onepicture[w=4.6cm]{figs/05.png}
\fourchoices[answer=A]
{等于零}
{不为零，方向向右}
{不为零，方向向左}
{不为零，$v_{0}$ 较大时方向向左，$v_{0}$ 较小时方向向右}

%% answer: A
%% memo:
\memoanswer{
斜劈和物块都平衡，对斜劈和物块整体受力分析知地面对斜劈的摩擦力为零，选 A。
}

\item
%% number: 6
%% paperName: 2011年海南高考第 6 题 3 分
%% typeId: 1
%% type: 单选题
%% chapter: 电磁感应
%% point: 电磁感应中图象
%% method: 无
%% score: 3
%% degree: 450
%% duplicateId: 0
%% body:
如图，$EOF$ 和 $E'O'F'$ 为空间一匀强磁场的边界，其中 $EO\parallel E'O'$，$FO\parallel F'O'$，且 $EO\perp OF$；$OO'$ 为 $\angle EOF$ 的角平分线，$OO'$ 间的距离为 $l$；磁场方向垂直于纸面向里。一边长为 $l$ 的正方形导线框沿 $OO'$ 方向匀速通过磁场，$t=0$ 时刻恰好位于图示位置。规定导线框中感应电流沿逆时针方向时为正，则感应电流 $i$ 与时间 $t$ 的关系图线可能正确的是\xzanswer{B}
\onepicture[w=4.8cm]{figs/06.png}
\fourchoices[answer=B, ispicture=true, h=1.5cm]
{figs/06a.png}
{figs/06b.png}
{figs/06c.png}
{figs/06d.png}

%% answer: B
%% memo:
\memoanswer{
从图示位置到左边框运动至 $O'$ 点，感应电流均匀增大为正；左边框再运动至 $OO'$ 中点过程，感应电流为正不变；左边框由 $OO'$ 中点再运动至 $O$ 过程感应电流先减小后方向增大；以后右边框再运动至 $OO'$ 中点过程，感应电流为负不变；右边框再运动至 $O$ 过程感应电流减小至 0，图 B 正确。
}

\item
%% number: 7
%% paperName: 2011年海南高考第 7 题 4 分
%% typeId: 2
%% type: 多选题
%% chapter: 其他
%% point: 物理学史
%% method: 无
%% score: 4
%% degree: 650
%% duplicateId: 0
%% body:
自然界的电、热和磁等现象都是相互联系的，很多物理学家为寻找它们之间的联系做出了贡献。下列说法正确的是\xzanswer{ACD}
\fourchoices[answer=ACD]
{奥斯特发现了电流的磁效应，揭示了电现象和磁现象之间的联系}
{欧姆发现了欧姆定律，说明了热现象和电现象之间存在联系}
{法拉第发现了电磁感应现象，揭示了磁现象和电现象之间的联系}
{焦耳发现了电流的热效应，定量得出了电能和热能之间的转换关系}

%% answer: ACD
%% memo:
\memoanswer{
考察物理学的发展史，选 ACD。
}

\item
%% number: 8
%% paperName: 2011年海南高考第 8 题 4 分
%% typeId: 2
%% type: 多选题
%% chapter: 直线运动
%% point: 运动图象
%% method: 无
%% score: 4
%% degree: 450
%% duplicateId: 0
%% body:
一物体自 $t=0$ 时开始做直线运动，其速度图线如图所示。下列选项正确的是\xzanswer{BC}
\onepicture[w=4.2cm]{figs/08.png}
\fourchoices[answer=BC]
{在 $0\sim6\Us$ 内，物体离出发点最远为 $30\Um$}
{在 $0\sim6\Us$ 内，物体经过的路程为 $40\Um$}
{在 $0\sim4\Us$ 内，物体的平均速率为 $7.5\Ums$}
{在 $5\sim6\Us$ 内，物体所受的合外力做负功}

%% answer: BC
%% memo:
\memoanswer{
在 $0\sim5\Us$ 内，物体向正向运动，$5\sim6\Us$ 内物体向负向运动，故 $5\Us$ 末离出发点最远，A 错；由面积法求出 $0\sim5\Us$ 内的位移 $s_{1}=35\Um$，$5\sim6\Us$ 内的位移 $s_{2}=-5\Um$，总路程为 $40\Um$，B 对；由面积法求出 $0\sim4\Us$ 内的位移 $s=30\Um$，平均速率为 $v=\frac{s}{t}=7.5\Ums$，C 对；由图像知 $5\sim6\Us$ 过程物体加速，合力和位移同向，合力做正功，D 错。
}

\item
%% number: 9
%% paperName: 2011年海南高考第 9 题 4 分
%% typeId: 2
%% type: 多选题
%% chapter: 功和能
%% point: 功率
%% method: 无
%% score: 4
%% degree: 450
%% duplicateId: 0
%% body:
一质量为 $1\Ukg$ 的质点静止于光滑水平面上，从 $t=0$ 时起，第 1 秒内受到 $2\UN$ 的水平外力作用，第 2 秒内受到同方向的 $1\UN$ 的外力作用。下列判断正确的是\xzanswer{CD}
\fourchoices[answer=CD]
{$0\sim2\Us$ 内外力的平均功率是 $\frac{9}{4}\UW$}
{第 2 秒内外力所做的功是 $\frac{5}{4}\UJ$}
{第 2 秒末外力的瞬时功率最大}
{第 1 秒内与第 2 秒内质点动能增加量的比值是 $\frac{4}{5}$}

%% answer: CD
%% memo:
\memoanswer{
由动量定理求出 $1\Us$ 末、$2\Us$ 末速度分别为 $v_{1}=2\Ums$、$v_{2}=3\Ums$，故合力做功为 $W=\frac{1}{2}mv^{2}=4.5\UJ$，功率为 $p=\frac{W}{t}=\frac{4.5}{3}\UW=1.5\UW$。$1\Us$ 末、$2\Us$ 末功率分别为 $4\UW$、$3\UW$；第 1 秒内与第 2 秒内动能增加量分别为 $\frac{1}{2}mv_{1}^{2}=2\UJ$、$\frac{1}{2}mv_{2}^{2}-\frac{1}{2}mv_{1}^{2}=2.5\UJ$，比值 $4:5$。
}

\item
%% number: 10
%% paperName: 2011年海南高考第 10 题 4 分
%% typeId: 2
%% type: 多选题
%% chapter: 磁场
%% point: 带电粒子在磁场中的运动
%% method: 无
%% score: 4
%% degree: 450
%% duplicateId: 0
%% body:
空间存在方向垂直于纸面向里的匀强磁场，图中的正方形为其边界。一细束由两种粒子组成的粒子流沿垂直于磁场的方向从 O 点入射。这两种粒子带同种电荷，它们的电荷量、质量均不同，但其比荷相同，且都包含不同速率的粒子。不计重力。下列说法正确的是\xzanswer{BD}
\onepicture[w=3.0cm]{figs/10.png}
\fourchoices[answer=BD]
{入射速度不同的粒子在磁场中的运动时间一定不同}
{入射速度相同的粒子在磁场中的运动轨迹一定相同}
{在磁场中运动时间相同的粒子，其运动轨迹一定相同}
{在磁场中运动时间越长的粒子，其轨迹所对的圆心角一定越大}

%% answer: BD
%% memo:
\memoanswer{
在磁场中半径 $r=\frac{mv}{qB}$，运动时间 $t=\frac{\theta m}{qB}$（$\theta$ 为转过的圆心角），故 BD 正确。当粒子从 O 点所在的边上射出时，轨迹可以不同，但圆心角相同为 $180^\circ$，因而 AC 错。
}

\item
%% number: 11
%% paperName: 2011年海南高考第 11 题 4 分
%% typeId: 3
%% type: 填空题
%% chapter: 交流电
%% point: 变压器
%% method: 无
%% score: 4
%% degree: 650
%% duplicateId: 0
%% body:
如图，理想变压器原线圈与一 $10\UV$ 的交流电源相连，副线圈并联两个小灯泡 a 和 b，小灯泡 a 的额定功率为 $0.3\UW$，正常发光时电阻为 $30\UO$，已知两灯泡均正常发光，流过原线圈的电流为 $0.09\UA$，可计算出原、副线圈的匝数比为\tkanswer[1.8cm]{$10:3$}，流过灯泡 b 的电流为\tkanswer[1.8cm]{$0.2$}$\UA$。
\onepicture[w=4.6cm, align=right]{figs/11.png}

%% answer: 10:3，0.2
%% memo:
\memoanswer{
副线圈电压 $U_{2}=U_{a}=\sqrt{P_{a}R_{a}}=3\UV$，$\frac{n_{1}}{n_{2}}=\frac{U_{1}}{U_{2}}=\frac{10}{3}$，$U_{b}=U_{a}=3\UV$。由能量守恒得 $U_{1}I_{1}=P_{a}+U_{b}I_{b}$，代入数据得 $I_{b}=0.2\UA$。
}

\item
%% number: 12
%% paperName: 2011年海南高考第 12 题 4 分
%% typeId: 3
%% type: 填空题
%% chapter: 曲线运动
%% point: 人造地球卫星
%% method: 无
%% score: 4
%% degree: 850
%% duplicateId: 0
%% body:
2011 年 4 月 10 日，我国成功发射第 8 颗北斗导航卫星，建成以后北斗导航卫星系统将包含多颗地球同步卫星，这有助于减少我国对 GPS 导航系统的依赖。GPS 由运行周期为 12 小时的卫星群组成，设北斗星的同步卫星和 GPS 导航的轨道半径分别为 $R_{1}$ 和 $R_{2}$，向心加速度分别为 $a_{1}$ 和 $a_{2}$，则 $R_{1}:R_{2}=$\tkanswer[2cm]{$\sqrt[3]{4}$}，$a_{1}:a_{2}=$\tkanswer[2cm]{$\frac{1}{2\sqrt[3]{2}}$}。（可用根式表示）

%% answer: ∛4，1/(2∛2)
%% memo:
\memoanswer{
$\frac{T_{1}}{T_{2}}=2$，由 $\frac{GMm}{R^{2}}=m\frac{4\pi^{2}}{T^{2}}R=ma$ 得 $R=\sqrt[3]{\frac{GMT^{2}}{4\pi^{2}}}$，$a=\frac{GM}{R^{2}}$，因而 $\frac{R_{1}}{R_{2}}=\left(\frac{T_{1}}{T_{2}}\right)^{\frac{2}{3}}=\sqrt[3]{4}$，$\frac{a_{1}}{a_{2}}=\left(\frac{R_{1}}{R_{2}}\right)^{-2}=\frac{1}{2\sqrt[3]{2}}$。
}

\item
%% number: 13
%% paperName: 2011年海南高考第 13 题 6 分
%% typeId: 4
%% type: 实验题
%% chapter: 电路
%% point: 电路实验
%% method: 无
%% score: 6
%% degree: 650
%% duplicateId: 0
%% body:
图\ref{2011海南13a}是改装并校准电流表的电路图，已知表头的量程为 $I_{g}=600\UuA$，内阻为 $R_{g}$，mA 是标准电流表，要求改装后的电流表量程为 $I=60\UmA$。完成下列填空。
\begin{enumerate}
	\item
	图\ref{2011海南13a}中分流电阻 $R_{p}$ 的阻值应为\tkanswer[2cm]{$\frac{I_{g}R_{g}}{I-I_{g}}$}（用 $I_{g}$、$R_{g}$ 和 $I$ 表示）。
	\item
	在电表改装完成后的某次校准测量中，mA 表的示数如图\ref{2011海南13b}所示，由此读出流过 mA 电流表的电流为\tkanswer[1.8cm]{$49.5$}$\UmA$。此时流过分流电阻 $R_{p}$ 的电流为\tkanswer[1.8cm]{$49.0$}$\UmA$（保留 1 位小数）。
\end{enumerate}
\twopicture[labelA=2011海南13a, labelB=2011海南13b, numA=1, numB=2, numstyle=arabic, h=3.2cm, align=right]
{figs/13a.png}
{figs/13b.png}

%% answer: 
\jdanswer{
\begin{enumerate}
	\item
	$\frac{I_{g}R_{g}}{I-I_{g}}$
	\item
	$49.5$，$49.0$
\end{enumerate}
}
%% memo:
\memoanswer{
（1）电流表改装，$R_{p}=\frac{1}{\frac{I}{I_{g}}-1}R_{g}=\frac{1}{\frac{60}{600\times10^{-3}}-1}R_{g}=\frac{R_{g}}{99}$。

（2）由图示可得，电流表读数为 $49.5\UmA$，实际流经表头的电流为 $500\UuA$ 即 $0.5\UmA$，此时流过分流电阻 $R_{p}$ 的电流为 $(49.5-0.5)\UmA=49.0\UmA$。
}

\item
%% number: 14
%% paperName: 2011年海南高考第 14 题 9 分
%% typeId: 4
%% type: 实验题
%% chapter: 功和能
%% point: DIS研究机械能守恒定律
%% method: 无
%% score: 9
%% degree: 250
%% duplicateId: 0
%% body:
现要通过实验验证机械能守恒定律。实验装置如图\ref{2011海南14a}所示：水平桌面上固定一倾斜的气垫导轨；导轨上 A 点处有一带长方形遮光片的滑块，其总质量为 $M$，左端由跨过轻质光滑定滑轮的细绳与一质量为 $m$ 的砝码相连；遮光片两条长边与导轨垂直；导轨上 B 点有一光电门，可以测量遮光片经过光电门时的挡光时间 $t$。用 $d$ 表示 A 点到导轨低端 C 点的距离，$h$ 表示 A 与 C 的高度差，$b$ 表示遮光片的宽度，$s$ 表示 A、B 两点的距离，将遮光片通过光电门的平均速度看作滑块通过 B 点时的瞬时速度。用 $g$ 表示重力加速度。完成下列填空和作图。
\onepicture[num=1, numstyle=arabic, label=2011海南14a, h=3.4cm, align=right]{figs/14a.png}
\begin{enumerate}
	\item
	若将滑块自 A 点由静止释放，则在滑块从 A 运动至 B 的过程中，滑块、遮光片与砝码组成的系统重力势能的减小量可表示为\tkanswer[2.4cm]{$(\frac{h}{d}M-m)gs$}。动能的增加量可表示为\tkanswer[2.4cm]{$\frac{1}{2}(M+m)\frac{b^{2}}{t^{2}}$}。若在运动过程中机械能守恒，$\frac{1}{t^{2}}$ 与 $s$ 的关系式为 $\frac{1}{t^{2}}=$\tkanswer[2.8cm]{$\frac{2(hM-dm)g}{(M+m)db^{2}}s$}。
	\item
	多次改变光电门的位置，每次均令滑块自同一点（A 点）下滑，测量相应的 $s$ 与 $t$ 值，结果如下表所示：
	\begin{center}
	\begin{tabular}{|c|c|c|c|c|c|}
	\hline
	 & 1 & 2 & 3 & 4 & 5 \\ \hline
	$s$（$\Um$） & 0.600 & 0.800 & 1.000 & 1.200 & 1.400 \\ \hline
	$t$（$\UmsT$） & 8.22 & 7.17 & 6.44 & 5.85 & 5.43 \\ \hline
	$1/t^{2}$（$10^{4}\Usnq$） & 1.48 & 1.95 & 2.41 & 2.92 & 3.39 \\ \hline
	\end{tabular}
	\end{center}
	以 $s$ 为横坐标，$\frac{1}{t^{2}}$ 为纵坐标，在答题卡上对应图\ref{2011海南14b}位置的坐标纸中描出第 1 和第 5 个数据点；根据 5 个数据点作直线，求得该直线的斜率 $k=$\tkanswer[1.6cm]{$2.40$}$\times10^{4}\Umn\cdot\Usnq$（保留 3 位有效数字）。
	\onepicture[num=2, numstyle=arabic, label=2011海南14b, w=6.0cm, align=right]{figs/14b.png}
	由测得的 $h$、$d$、$b$、$M$ 和 $m$ 数值可以计算出 $\frac{1}{t^{2}}-s$ 直线的斜率 $k_{0}$，将 $k$ 和 $k_{0}$ 进行比较，若其差值在试验允许的范围内，则可认为此试验验证了机械能守恒定律。
\end{enumerate}

%% answer: 
\jdanswer{
\begin{enumerate}
	\item
	$(\frac{h}{d}M-m)gs$；$\frac{1}{2}(M+m)\frac{b^{2}}{t^{2}}$；$\frac{2(hM-dm)g}{(M+m)db^{2}}s$
	\item
	$2.40$
\end{enumerate}
}
%% memo:
\memoanswer{
（1）重力势能的减少量 $\Delta E_{p}=Mg\frac{h}{d}s-mgs=(\frac{h}{d}M-m)gs$；到达 B 点的速度为 $v_{B}=\frac{b}{t}$，动能的增加量 $\Delta E_{k}=\frac{1}{2}(M+m)v_{B}^{2}=\frac{1}{2}(M+m)\frac{b^{2}}{t^{2}}$。若在运动过程中机械能守恒，$\Delta E_{p}=\Delta E_{k}$ 可得 $\frac{1}{t^{2}}=\frac{2(hM-dm)g}{(M+m)db^{2}}s$。

（2）利用图像可求得斜率为 $2.40\times10^{4}\Umn\cdot\Usnq$。利用实验数据分别进行逐差法得直线斜率为 $2.399\times10^{4}\Umn\cdot\Usnq$，保留 3 位有效数字为 $2.40\times10^{4}\Umn\cdot\Usnq$。

\onepicture[w=5.5cm]{figs/14_an.png}
}

\item
%% number: 15
%% paperName: 2011年海南高考第 15 题 8 分
%% typeId: 6
%% type: 计算题
%% chapter: 曲线运动
%% point: 平抛运动
%% method: 无
%% score: 8
%% degree: 450
%% duplicateId: 0
%% body:
如图，水平地面上有一个坑，其竖直截面为半圆。ab 为沿水平方向的直径。若在 a 点以初速度 $v_{0}$ 沿 ab 方向抛出一小球，小球会击中坑壁上的 c 点。已知 c 点与水平地面的距离为圆半径的一半，求圆的半径。
\onepicture[w=4.0cm, align=right]{figs/15.png}

%% answer: 
\jdanswer{
$r=\frac{4(7\pm4\sqrt{3})v_{0}^{2}}{g}$
}
%% memo:
\memoanswer{
设圆半径为 $r$，小球做平抛运动，则：
$x=v_{0}t$ ①
$y=0.5r=\frac{1}{2}gt^{2}$ ②
过 c 点作 cd$\perp$ab 交 ab 于 d 点，$\triangle acd\sim\triangle cbd$ 可得 $cd^{2}=ad\cdot db$，即：
$\left(\frac{r}{2}\right)^{2}=x(2r-x)$ ③
由①②③式得 $r=\frac{4(7\pm4\sqrt{3})v_{0}^{2}}{g}$。
}

\item
%% number: 16
%% paperName: 2011年海南高考第 16 题 11 分
%% typeId: 6
%% type: 计算题
%% chapter: 电磁感应
%% point: 双杆问题
%% method: 无
%% score: 11
%% degree: 450
%% duplicateId: 0
%% body:
如图，ab 和 cd 是两条竖直放置的长直光滑金属导轨，$MN$ 和 $M'N'$ 是两根用细线连接的金属杆，其质量分别为 $m$ 和 $2m$。竖直向上的外力 $F$ 作用在杆 $MN$ 上，使两杆水平静止，并刚好与导轨接触；两杆的总电阻为 $R$，导轨间距为 $l$。整个装置处在磁感应强度为 $B$ 的匀强磁场中，磁场方向与导轨所在平面垂直。导轨电阻可忽略，重力加速度为 $g$。在 $t=0$ 时刻将细线烧断，保持 $F$ 不变，金属杆和导轨始终接触良好。求：
\begin{enumerate}
	\item
	细线烧断后，任意时刻两杆运动的速度之比；
	\item
	两杆分别达到的最大速度。
\end{enumerate}
\onepicture[w=3.4cm, align=right]{figs/16.png}

%% answer: 
\jdanswer{
\begin{enumerate}
	\item
	$2:1$
	\item
	$v_{1}=\frac{4mgR}{3B^{2}l^{2}}$，$v_{2}=\frac{2mgR}{3B^{2}l^{2}}$
\end{enumerate}
}
%% memo:
\memoanswer{
设某时刻 $MN$ 和 $M'N'$ 的速度分别为 $v_{1}$、$v_{2}$。

（1）$MN$ 和 $M'N'$ 动量守恒：$mv_{1}-2mv_{2}=0$，求出 $\frac{v_{1}}{v_{2}}=2$ ①。

（2）当 $MN$ 和 $M'N'$ 的加速度为零时，速度最大。对 $M'N'$ 受力平衡：$BIl=2mg$ ②；$I=\frac{E}{R}$ ③；$E=Blv_{1}+Blv_{2}$ ④。由①②③④式得 $v_{1}=\frac{4mgR}{3B^{2}l^{2}}$，$v_{2}=\frac{2mgR}{3B^{2}l^{2}}$。
}

\item
%% number: 17
%% paperName: 2011年海南高考第 17 题 8 分
%% typeId: 6
%% type: 计算题
%% chapter: 气体性质
%% point: 玻意耳定律
%% method: 无
%% score: 8
%% degree: 650
%% duplicateId: 0
%% body:
【模块 3-3 试题】（12 分）
\begin{enumerate}
	\item
	（4 分）关于空气湿度，下列说法正确的是\tkanswer[1.6cm]{AC}。（填入正确选项前的字母。选对 1 个给 2 分，选对 2 个给 4 分；选错 1 个扣 2 分，最低得 0 分）
	\fourchoices[answer=AC]
	{当人们感到潮湿时，空气的绝对湿度一定较大}
	{当人们感到干燥时，空气的相对湿度一定较小}
	{空气的绝对湿度用空气中所含水蒸汽的压强表示}
	{空气的相对湿度定义为水的饱和蒸汽压与相同温度时空气中所含水蒸气的压强之比}
	\item
	（8 分）如图，容积为 $V_{1}$ 的容器内充有压缩空气。容器与水银压强计相连，压强计左右两管下部由软胶管相连。气阀关闭时，两管中水银面等高，左管中水银面上方到气阀之间空气的体积为 $V_{2}$。打开气阀，左管中水银下降；缓慢地向上提右管，使左管中水银面回到原来高度，此时右管与左管中水银面的高度差为 $h$。已知水银的密度为 $\rho$，大气压强为 $p_{0}$，重力加速度为 $g$；空气可视为理想气体，其温度不变。求气阀打开前容器中压缩空气的压强 $p_{1}$。
	\onepicture[w=6.5cm, align=right]{figs/17.png}
\end{enumerate}

%% answer: 
\jdanswer{
\begin{enumerate}
	\item
	AC
	\item
	$p_{1}=\frac{\rho gh(V_{1}+V_{2})+p_{0}V_{1}}{V_{1}}$
\end{enumerate}
}
%% memo:
\memoanswer{
（1）人们感到潮湿时，空气中实际水汽含量大，即空气中所含水蒸气的压强大，空气的绝对湿度大，A、C 正确；空气的相对湿度定义为水蒸气的实际压强与相同温度时水的饱和蒸汽压之比，人们感到干燥时，相对湿度较大，B、D 错误。

（2）由玻意耳定律得 $p_{1}V_{1}+p_{0}V_{2}=(p_{0}+\rho gh)(V_{1}+V_{2})$，$p_{1}=\frac{\rho gh(V_{1}+V_{2})+p_{0}V_{1}}{V_{1}}$。
}

\item
%% number: 18
%% paperName: 2011年海南高考第 18 题 8 分
%% typeId: 6
%% type: 计算题
%% chapter: 光学
%% point: 光的折射
%% method: 无
%% score: 8
%% degree: 650
%% duplicateId: 0
%% body:
【模块 3-4 试题】（12 分）
\begin{enumerate}
	\item
	（4 分）一列简谐横波在 $t=0$ 时的波形图如图所示。介质中 $x=2\Um$ 处的质点 P 沿 $y$ 轴方向做简谐运动的表达式为 $y=10\sin(5\pi t)\Ucm$。关于这列简谐波，下列说法正确的是\tkanswer[1.6cm]{CD}。（填入正确选项前的字母。选对 1 个给 2 分，选对 2 个给 4 分；选错 1 个扣 2 分，最低得 0 分）
	\onepicture[w=5.0cm]{figs/18.png}
	\fourchoices[answer=CD]
	{周期为 $4.0\Us$}
	{振幅为 $20\Ucm$}
	{传播方向沿 $x$ 轴正向}
	{传播速度为 $10\Ums$}
	\item
	（8 分）一赛艇停在平静的水面上，赛艇前端有一标记 P 离水面的高度为 $h_{1}=0.6\Um$，尾部下端 Q 略高于水面；赛艇正前方离赛艇前端 $s_{1}=0.8\Um$ 处有一浮标，示意如图。一潜水员在浮标前方 $s_{2}=3.0\Um$ 处下潜到深度为 $h_{2}=4.0\Um$ 时，看到标记刚好被浮标挡住，此处看不到船尾端 Q；继续下潜 $\Delta h=4.0\Um$，恰好能看见 Q。求：
	\begin{enumerate}[label = (\roman*)]
		\item
		水的折射率 $n$；
		\item
		赛艇的长度 $l$。（可用根式表示）
	\end{enumerate}
	\onepicture[w=4.5cm, align=right]{figs/18b.png}
\end{enumerate}

%% answer: 
\jdanswer{
\begin{enumerate}
	\item
	CD
	\item
	（i）$n=\frac{4}{3}$；（ii）$l=\left(\frac{8\sqrt{7}}{3}-3.8\right)\Um\approx3.3\Um$。
\end{enumerate}
}
%% memo:
\memoanswer{
（1）$\omega=5\pi$，周期为 $T=\frac{2\pi}{\omega}=0.4\Us$；由波的图像得振幅 $A=10\Ucm$、波长 $\lambda=4\Um$，故波速为 $v=\frac{\lambda}{T}=10\Ums$；P 点在 $t=0$ 时振动方向为正 $y$ 方向，波向正 $x$ 方向传播。

（2）（i）设过 P 点的光线恰好被浮标挡住时，入射角、折射角分别为 $\alpha$、$\beta$，则 $\sin\alpha=\frac{s_{1}}{\sqrt{s_{1}^{2}+h_{1}^{2}}}$ ①，$\sin\beta=\frac{s_{2}}{\sqrt{s_{2}^{2}+h_{2}^{2}}}$ ②，$n=\frac{\sin\alpha}{\sin\beta}$ ③，由①②③式得 $n=\frac{4}{3}$。

（ii）潜水员和 Q 点连线与水平方向夹角刚好为临界角 C，则 $\sin C=\frac{1}{n}=\frac{3}{4}$ ④，$\tan C=\frac{h_{2}+\Delta h}{s_{1}+s_{2}+l}$ ⑤，由④⑤式得 $l=\left(\frac{8\sqrt{7}}{3}-3.8\right)\Um\approx3.3\Um$。
}

\item
%% number: 19
%% paperName: 2011年海南高考第 19 题 8 分
%% typeId: 6
%% type: 计算题
%% chapter: 运动和力
%% point: 动量和能量
%% method: 无
%% score: 8
%% degree: 450
%% duplicateId: 0
%% body:
【模块 3-5 试题】（12 分）
\begin{enumerate}
	\item
	（4 分）2011 年 3 月 11 日，日本发生九级大地震，造成福岛核电站的核泄漏事故。在泄漏的污染物中含有 $^{131}$I 和 $^{137}$Cs 两种放射性核素，它们通过一系列衰变产生对人体有危害的辐射。在下列四个式子中，有两个能分别反映 $^{131}$I 和 $^{137}$Cs 的衰变过程，它们分别是\tkanswer[1.2cm]{B}和\tkanswer[1.2cm]{C}（填入正确选项前的字母）。$^{131}$I 和 $^{137}$Cs 原子核中的中子数分别是\tkanswer[1.2cm]{78}和\tkanswer[1.2cm]{82}。
	\fourchoices[answer=BC]
	{$\ce{X_{1} -> ^{137}_{56}Ba + ^{1}_{0}n}$}
	{$\ce{X_{2} -> ^{131}_{54}Xe + ^{0}_{-1}e}$}
	{$\ce{X_{3} -> ^{137}_{56}Ba + ^{0}_{-1}e}$}
	{$\ce{X_{4} -> ^{131}_{54}Xe + ^{1}_{1}p}$}
	\item
	（8 分）一质量为 $2m$ 的物体 P 静止于光滑水平地面上，其截面如图所示。图中 ab 为粗糙的水平面，长度为 $L$；bc 为一光滑斜面，斜面和水平面通过与 ab 和 bc 均相切的长度可忽略的光滑圆弧连接。现有一质量为 $m$ 的木块以大小为 $v_{0}$ 的水平初速度从 a 点向左运动，在斜面上上升的最大高度为 $h$，返回后在到达 a 点前与物体 P 相对静止。重力加速度为 $g$。求：
	\begin{enumerate}[label = (\roman*)]
		\item
		木块在 ab 段受到的摩擦力 $f$；
		\item
		木块最后距 a 点的距离 $s$。
	\end{enumerate}
	\onepicture[w=6.5cm, align=right]{figs/19.png}
\end{enumerate}

%% answer: 
\jdanswer{
\begin{enumerate}
	\item
	B、C、78、82
	\item
	（i）$f=\frac{m(v_{0}^{2}-3gh)}{3L}$；（ii）$s=\frac{v_{0}^{2}-6gh}{v_{0}^{2}-3gh}L$。
\end{enumerate}
}
%% memo:
\memoanswer{
（1）由质量数和核电荷数守恒可以得出正确选项 B、C。$^{131}$I 和 $^{137}$Cs 原子核中的中子数分别为 78 和 82。

（2）（i）设木块和物体 P 的共同速度为 $v$，两物体从开始到第一次达到共同速度的过程由动量和能量守恒得 $mv_{0}=(m+2m)v$ ①，$\frac{1}{2}mv_{0}^{2}=\frac{1}{2}(m+2m)v^{2}+mgh+fL$ ②，由①②式得 $f=\frac{m(v_{0}^{2}-3gh)}{3L}$ ③。

（ii）木块返回与物体 P 第二次达到共同速度与第一次相同（动量守恒），全过程能量守恒得 $\frac{1}{2}mv_{0}^{2}=\frac{1}{2}(m+2m)v^{2}+f(2L-s)$ ④，由②③④式得 $s=\frac{v_{0}^{2}-6gh}{v_{0}^{2}-3gh}L$。
}

\end{enumerate}

\end{document}
